\section{The boundary-selected infinite-volume state}\label{sec:topology}

The gap just proved supplies the hypothesis of Tasaki's topological-index
theorem. We formulate the consequence for subsequential local limits,
so that no convergence or uniqueness of the infinite-volume state is
presupposed.

A local observable is a matrix acting on finitely many sites of
$\mathbb Z$, extended by the identity on the other sites. A state on
these observables is a normalized positive linear functional. Write
$\phi_L$ for a normalized ground vector of $H_L(3/5)$ and set
$\omega_L(A)=\langle\phi_L,A\phi_L\rangle$ whenever the support of
$A$ lies in $[-L,L]$. For $L\ge960$ the vector is unique up to phase
by Theorem~\ref{thm:main}, so this functional is well defined.
Compactness of the density matrices on each finite interval, followed
by a diagonal subsequence over the intervals, gives integers
$L_k\to\infty$ for which
\begin{equation}\label{top:limit}
 \omega(A)=\lim_{k\to\infty}\omega_{L_k}(A)
\end{equation}
exists for every local observable. The limits are consistent on nested
intervals and therefore define a state. We call any such state
\emph{boundary-selected}.

The bulk Hamiltonian is the formal sum
$H_{\rm bulk}=\sum_{j\in\mathbb Z}S_j\cdot S_{j+1}$.
For a local $A$, its commutator $[H_{\rm bulk},A]$ is the finite sum
of the commutators with bonds meeting the support of $A$.
To specify the index, for $\varepsilon>0$ define the local twist
\begin{equation}\label{top:twist}
 U_\varepsilon=
 \exp\left[-i\sum_{j\in\mathbb Z\cap[-\pi/\varepsilon,\pi/\varepsilon]}
                 (\pi+\varepsilon j)S_j^z\right].
\end{equation}
The rotations by $0$ and $2\pi$ outside this interval act trivially
for spin one. For the symmetric gapped states considered below,
Tasaki's theorem guarantees the existence of the index
$\operatorname{Ind}(\omega)=\lim_{\varepsilon\downarrow0}
\omega(U_\varepsilon)\in\{-1,1\}$; we use precisely the twist convention in
\cite[Equations (10)--(13)]{Tasaki2025}.

\begin{corollary}\label{cor:topology}
Every boundary-selected state \eqref{top:limit} is invariant under
uniform rotations about the $z$ axis and under reflection $j\mapsto-j$.
With $\gamma_*=(\log10)/392$, it satisfies
\begin{equation}\label{top:local-gap}
 \omega\bigl(A^*[H_{\rm bulk},A]\bigr)
 \ge\gamma_*\bigl(\omega(A^*A)-|\omega(A)|^2\bigr)
 \qquad\text{for every local }A.
\end{equation}
In particular it is a locally-unique gapped ground state in the sense
of \cite[Definition 11]{Tasaki2025}, and its Tasaki index is $-1$.
\end{corollary}

\begin{proof}
Both finite-volume symmetries commute with $H_L(3/5)$. Simplicity of
the ground eigenspace implies invariance of $\omega_L$ under them;
passing to \eqref{top:limit} proves the asserted invariances.

For $L\ge960$, the finite spectral theorem and Theorem~\ref{thm:main}
give
\[
 \langle A\phi_L,(H_L(3/5)-E_0(L,3/5))A\phi_L\rangle
 \ge\gamma_*\bigl(\omega_L(A^*A)-|\omega_L(A)|^2\bigr).
\]
The expression on the left is
$\omega_L(A^*[H_L(3/5),A])$. Once the support of $A$ and its adjacent
sites lie strictly inside $[-L,L]$, this commutator equals
$[H_{\rm bulk},A]$: the endpoint fields and all omitted bonds commute
with $A$. Taking the limit proves \eqref{top:local-gap}. Its right
side is nonnegative by the Cauchy--Schwarz inequality for a state;
for $\omega(A)=0$ it is the local gap inequality. The strict
finite-volume lower bound has consequently yielded the non-strict
lower bound $\gamma_*$ in the limit.

Finally the Hamiltonians, spin normalization and endpoint-field signs
are exactly those of \cite[Equation (2)]{Tasaki2025}. Theorem~\ref{thm:main}
verifies that paper's Assumption 2 with the single field $h=3/5$,
$L_0=960$ and $\gamma=\gamma_*$. Its Theorem 3, with the subsequential
interpretation specified in its End Matter, Equation (27), therefore
applies to every state \eqref{top:limit}. To make the sign and limit
order explicit, Tasaki's finite-chain twist interpolation begins at
the uniform $\pi$ rotation. Its expectation is $-1$, because the
finite ground state has total $S^z=1$ by that paper's Lemma 1.
Reflection makes the twist expectation real, and the uniform gap
prevents it from crossing zero while the twist parameter is sufficiently
small. More precisely,
\cite[Lemmas 6--8 and Equations (17)--(18)]{Tasaki2025} gives, for
every sufficiently small fixed $\varepsilon>0$ and all sufficiently
large $L$,
\[
 -1\le\omega_L(U_\varepsilon)
 \le-\sqrt{1-\frac{4(\pi+\varepsilon)\varepsilon}{\gamma_*}}.
\]
Here $L\ge\pi/\varepsilon$ ensures that the finite-chain twist equals
the local operator \eqref{top:twist}. Passing first to
\eqref{top:limit} and then letting $\varepsilon\downarrow0$ gives
$\operatorname{Ind}(\omega)=-1$.
\end{proof}

The corollary evaluates the twist index for the selected states. It does
not require the full sequence $\omega_L$ to converge or identify these
states with every ground state of the infinite chain. In particular,
the additional global uniqueness hypothesis of the half-chain conclusion
in \cite[Corollary 4]{Tasaki2025} is not part of our result.
